\documentclass[11pt,a4paper]{article}
\usepackage[utf8]{inputenc}
\usepackage[T1]{fontenc}
\usepackage[margin=2.5cm]{geometry}
\usepackage{booktabs}
\usepackage{amsmath}
\usepackage{siunitx}
\usepackage{multirow}
\usepackage{float}
\usepackage[autostyle]{csquotes}
\usepackage[hidelinks]{hyperref}
\usepackage{graphicx}

\sisetup{detect-weight=true}

\title{\bfseries Marocco FL thesis --- week 12/10 update}
\author{}
\date{}

\begin{document}
\maketitle
\vspace{-3.2em}
\begin{center}
\emph{Client-Driven Local Work Assignment in Battery-Constrained Federated Learning}
\end{center}
\vspace{1em}

%======================================================================
\section*{Summary}

\begin{itemize}
  \item This week every physical parameter of the simulator was checked and,
        where possible, \textbf{tied to a published measurement}. Table~\ref{tab:overview}
        lists every value that changed, with its source.
  \item The three device tiers are now three \textbf{real phones} measured
        while training (Mandheling, MobiCom 2022): speed \emph{and} power come
        from the same measurement. The made-up power law
        $P(B)=\SI{3}{\watt}\,(B/128)^{0.3}$ is gone: the training power is
        constant per tier.
  \item Communication is now a \textbf{4G (LTE)} smartphone measured in 2021
        (Narayanan et al., SIGCOMM 2021) on Italian 4G speeds (Opensignal):
        0.98~J/MB down, 2.96~J/MB up, \SI{2.1}{\joule} per radio
        session. One round of model exchange costs \SI{8.5}{\joule} instead of
        \SI{29.2}{\joule}.
  \item Peukert exponent $n=1.05$ (fit on the datasheet capacities, was 1.15);
        recharge is now an \textbf{hourly} rate, $0.126$/h, derived from a
        study of 4035 users (Ferreira et al., Pervasive 2011).
  \item A \textbf{batch-size sweep} ($B$ from 8 to 256) shows that a smaller
        batch always gives more accuracy. SAGE-smart, which adapts the number
        of epochs, gets it for free in battery; FedAvg pays for it with
        depleted devices. SAGE-smart now chooses $B$ itself (16--64); the
        baselines stay at $B=128$.
  \item Making $B$ vary across clients broke the FedNova aggregation of
        SAGE-smart (it also extrapolated the BatchNorm statistics). Fixed.
  \item The full campaign with the new physics (33 runs) is running;
        results will be added in Section~\ref{sec:results}.
\end{itemize}

%======================================================================
\section{Overview: what changed}

\begin{table}[H]
\centering
\small
\caption{Every physical value that changed this week. ``Before'' is the
value used in the campaign of the 05/10 report.}
\label{tab:overview}
\begin{tabular}{p{3.4cm}p{3.2cm}p{3.6cm}p{4.2cm}}
\toprule
parameter & before & now & source \\
\midrule
tier speed (low/mid/high) & 4 / 8 / 16 GMAC/s, no source & 13.9 / 26.6 / 30.5 GMAC/s & Mandheling, Fig.~5 + Tab.~6 \\
training power & $3\,\mathrm{W}\,(B/128)^{0.3}$, same for all & 2.99 / 4.00 / 6.15 W, constant & Mandheling, Fig.~5--6 \\
power between rounds & $0.1 + w\,(3 - 0.1)$ W & $0.1 + w\,(P_\text{tier} - 0.1)$ W & consistency with the tier \\
fixed cost per batch & \SI{0.115}{\second}, undocumented & \SI{0.057}{\second} & Mandheling, Fig.~7 \\
training / forward MACs & 3, assumed & 3, now with a source & Mandheling Tab.~6; TrainDeeploy \\
Peukert exponent $n$ & 1.15 (lead-acid value) & 1.05 & fit on the datasheet capacities \\
4G download & 1.5 J/MB (supervisor) & 0.98 J/MB & Narayanan 2021 + Opensignal 2018 \\
4G upload & $1.5\times$ download, assumed & 2.96 J/MB & Narayanan 2021 + Opensignal 2018 \\
bandwidth & 2 MB/s both ways & 25.1 / 10 Mbit/s down / up & Opensignal 2018 \\
radio tail & 1.06 W $\times$ 11.6 s & 0.178 W $\times$ 10.2 s & Narayanan 2021, Tab.~2 and 7 \\
radio promotion & 1.21 W $\times$ 0.260 s & 1.21 W $\times$ 0.265 s & Huang 2012 / Narayanan 2021 \\
recharge & $p=0.005\,(1-\mathrm{SoC})$ per round & $p=1-e^{-0.126\,(1-\mathrm{SoC})\,\Delta t}$, $\Delta t$ in hours & Ferreira 2011 \\
SAGE-smart batch & $\{32, 64\}$ from the $P(B)$ optimum & chosen by the client in $\{16, 32, 64\}$ & batch sweep (Sec.~\ref{sec:batch}) \\
\bottomrule
\end{tabular}
\end{table}

Values that did \emph{not} change but now have an explicit source:
battery capacities 3500 / 4000 / 4500 mAh (a subset of the ``common smartphone
capacities'' used by ESCS), idle power \SI{0.1}{\watt} (ESCS), mean initial SoC
0.75 (Ferreira: average battery 74\,\% at 5 a.m., i.e.\ at the start of the
day), charge rate 1C (Ferreira: 1--2\,h for a full charge on 2010 phones,
faster today), reserve 0.2 (SAGE constraint 2e, ESCS \texttt{B\_LVL\_min} $=20\,\%$).
Still assumptions: the workload levels (0 / 0.1 / 0.3) and their Markov
chains, the standard deviation 0.15 of the initial SoC, the 1~KB control
message.

%======================================================================
\section{Device tiers: three real phones}

\paragraph{Source.} Mandheling (Xu et al., MobiCom 2022) trains CNNs in FP32 on
the CPU of three smartphones and reports, for a batch of 64 samples, both the
\emph{time} (Fig.~5) and the \emph{energy} (Fig.~6) of one training step. We
use TensorFlow Lite on VGG-11 / CIFAR-10, the same input size as ours.
Table~6 of the paper gives 914 MFLOP of training per sample, i.e.\
457~MMAC $= 3 	imes 152.6$~MMAC of forward pass: the same
training/forward ratio of 3 that we use.

\paragraph{How.}
\[
\text{speed} = \frac{64 \times 457\cdot 10^{6}\ \text{MAC}}{\text{time per batch}},
\qquad
\text{power} = \frac{\text{energy per batch}}{\text{time per batch}} .
\]

\begin{table}[H]
\centering
\caption{The three tiers. J/GMAC = energy per unit of work.}
\begin{tabular}{llrrrrr}
\toprule
tier & phone (SoC) & time [s] & energy [J] & GMAC/s & power [W] & J/GMAC \\
\midrule
low  & Redmi Note 9 Pro (Snapdragon 750G) & 2.11 & 6.3 & 13.9 & 2.99 & 0.215 \\
mid  & Xiaomi 10 (Snapdragon 865)         & 1.10 & 4.4 & 26.6 & 4.00 & 0.150 \\
high & Xiaomi 11 Pro (Snapdragon 888)     & 0.96 & 5.9 & 30.5 & 6.15 & 0.202 \\
\bottomrule
\end{tabular}
\end{table}

Cross-check: ResNet-18 on the same phones gives 16.6 / 22.0 / 26.2 GMAC/s at
2.98 / 4.11 / 6.16~W, i.e.\ the same power and a similar speed.

\paragraph{What it means.} The energy per unit of work is almost the same on
the three phones ($\approx 0.2$~J/GMAC): the fast phone finishes earlier but
draws more watts. Before, the low tier cost 0.75~J/GMAC and the high tier
0.19, a factor 4 with no source. For one epoch of ResNet-20 on 1666 samples
at $B=128$:

\begin{center}
\begin{tabular}{lrrrr}
\toprule
tier & before [J] & now [J] & now [s] & ratio \\
\midrule
low  & 156.5 & 45.8 & 15.3 & $3.4\times$ cheaper \\
mid  &  80.5 & 33.4 &  8.4 & $2.4\times$ cheaper \\
high &  42.5 & 45.4 &  7.4 & about the same \\
\bottomrule
\end{tabular}
\end{center}

\paragraph{Training power is constant per tier.} Mandheling gives one power
value per phone. During training the CPU runs at full load whatever the batch
size; $B$ changes how long the work lasts, not how many watts the phone draws.
The old law $P(B)\propto B^{0.3}$ had no source and has been removed.

\paragraph{Limits.} VGG-11 has 64--512 channels, ResNet-20 only 16--64: on a
CPU ResNet-20 probably reaches fewer GMAC/s. The measurements are per batch,
so they do not include thermal throttling over a long training. The three
phones are 2020--21 mid- and high-end; there is no budget phone.

%======================================================================
\section{Cost of a training step}

\paragraph{Training / forward ratio = 3.} Now with two sources: Mandheling
Table~6 ($914/2/152.6 = 3.0$) and the TrainDeeploy measurements on GAP9
(2.85--2.97).

\paragraph{Fixed cost per batch = \SI{0.057}{\second}.} The old value
(\SI{0.115}{\second}, ``measured on the GPU server and rescaled'') had no
written derivation and was not measured on a phone. It matters: it is the
only way $B$ enters the energy. Mandheling Fig.~7 plots the time of one
training batch against $B$ (4 to 128) on the Xiaomi 11 Pro, TFLite FP32. The
figure is vector graphics: we extracted the exact coordinates of the curves
from the PDF and calibrated the axis on the values at $B=64$ printed in
Fig.~5.

\begin{center}
\begin{tabular}{lrrrrrrr}
\toprule
$B$ & 4 & 8 & 16 & 32 & 64 & 96 & 128 \\
\midrule
VGG-16 [s]      & 0.12 & 0.18 & 0.31 & 0.54 & 1.39 & 1.57 & 2.07 \\
InceptionV3 [s] & 0.22 & 0.36 & 0.65 & 1.70 & 3.18 & 4.08 & 5.17 \\
\bottomrule
\end{tabular}
\end{center}

The time is almost proportional to $B$. The fixed part is what remains at the
smallest batch once the computation of 4 samples is removed:
\[
t_\text{ov} = t(4) - 4\,\frac{t(128)}{128}:\qquad
0.12 - 4\times 0.0162 = \SI{0.055}{\second},\qquad
0.22 - 4\times 0.0404 = \SI{0.058}{\second}.
\]
Two very different models give the same value; we use \SI{0.057}{\second}.
(MNN on the same phone gives 0.02--0.03~s.)

\paragraph{Effect.} With constant power, energy per sample $=P\,(c + t_\text{ov}/B)$,
with $c$ the compute time per sample. A smaller batch costs more per sample:

\begin{center}
\begin{tabular}{lrrrrrr}
\toprule
& \multicolumn{4}{c}{mJ per sample, now} & \multicolumn{2}{c}{$B=32$ vs $B=128$} \\
\cmidrule(lr){2-5}\cmidrule(lr){6-7}
tier & $B=16$ & $B=32$ & $B=64$ & $B=128$ & now & with \SI{0.115}{\second} \\
\midrule
low  & 36.8 & 31.5 & 28.8 & 27.5 & $+15\,\%$ & $+28\,\%$ \\
mid  & 32.5 & 25.4 & 21.9 & 20.1 & $+27\,\%$ & $+49\,\%$ \\
high & 46.4 & 35.5 & 30.0 & 27.3 & $+30\,\%$ & $+55\,\%$ \\
\bottomrule
\end{tabular}
\end{center}

%======================================================================
\section{Power between rounds}

Between rounds a phone draws
$P = P_\text{idle} + w\,(P_\text{train} - P_\text{idle})$, where
$P_\text{idle}=\SI{0.1}{\watt}$ (ESCS) and $w$ is the share of the CPU taken by
the user's other apps (0 screen off, 0.1 background, 0.3 interactive). $w$ is
a fraction of the capacity \emph{of that phone} (it already reduces the
training speed as $(1-w)\times$ speed), so $P_\text{train}$ is now the one of
the tier instead of \SI{3}{\watt} for everybody:

\begin{center}
\begin{tabular}{lrrr}
\toprule
& low & mid & high \\
\midrule
$w=0.1$ & 0.39 W & 0.49 W & 0.71 W \\
$w=0.3$ & 0.97 W & 1.27 W & 1.92 W \\
average $w=0.118$ & 0.44 W & 0.56 W & 0.82 W \\
full battery lasts (no FL) & 29 h & 26 h & 20 h \\
\bottomrule
\end{tabular}
\end{center}

A full battery lasting about a day with normal use is realistic, and it is
consistent with the one charge per day used for recharging (Sec.~\ref{sec:recharge}).

%======================================================================
\section{Battery: Peukert exponent}

The datasheet of the cell used in the datasheet world (``Accurate Battery
Model'') gives three capacities: 2.60 / 2.50 / 2.39~Ah at 0.2C / 0.5C / 1C.
Peukert's law gives $Q(I) = Q_R\,(I_R/I)^{\,n-1}$ with $I_R = 0.2$C (the rate at
which the nominal capacity is declared). A least-squares fit of
$\ln(Q/Q_R)$ against $\ln(I/I_R)$ through the rated point gives
\textbf{$n = 1.050$} (1.052 without the constraint).

The old 1.15 is a lead-acid value; for lithium cells the literature gives
1.01--1.10. Effect: with \SI{0.1}{\watt} of idle on 3500~mAh, a full Peukert
battery lasted 1.63 times the linear one; now 1.18 times. Limit: the fit
covers 0.2C--1C, while idle (about 0.01C) is an extrapolation.

The datasheet tables $V(\mathrm{SoC}, I)$ and $Q(I)$ are unchanged from the
05/10 report.

%======================================================================
\section{Communication: 4G}

\paragraph{Energy per MB.} Narayanan et al.\ (SIGCOMM 2021) measure with a
Monsoon power meter the power of a Galaxy S20 Ultra transferring data on
Verizon 4G/LTE. Power grows linearly with throughput $T$:
$P = c_1 T + c_2$.
\begin{itemize}
  \item $c_1$ from their Table~8: 14.55~mW per Mbit/s download,
        80.21~mW per Mbit/s upload (upload grows 2.2--5.9$\times$ faster
        than download, as they write).
  \item $c_2$ is not printed in the paper. We fitted it on their public
        dataset (GitHub \texttt{SIGCOMM21-5G/artifact}, \texttt{MN-Transfer}):
        2.72~W download, 2.89~W upload. The slopes of our fit (15.0 and 79.5)
        reproduce their Table~8, so we are using the right data.
\end{itemize}
Energy per MB $= 8\,(c_1 + c_2/T)$, so it depends on the throughput. We take
the average Italian 4G throughput from Opensignal (``State of Mobile Networks:
Italy'', May 2018, 120{,}106 users, 1.34 billion measurements, Jan--Mar 2018):

\begin{center}
\begin{tabular}{lrrrr|r}
\toprule
4G & Vodafone & TIM & 3 Italia & Wind & used \\
\midrule
download [Mbit/s] & 33.9 & 28.7 & 22.2 & 15.7 & 25.1 \\
upload [Mbit/s]   & 10.5 & \multicolumn{3}{c|}{``just below 10''} & 10 \\
\bottomrule
\end{tabular}
\end{center}

Result: \textbf{0.98~J/MB download, 2.96~J/MB upload} (ratio 3.0, inside the
2.2--5.9 of the paper). For comparison, the 1.5~J/MB we used before
corresponds to about 16~Mbit/s, i.e.\ the slowest operator. The same
throughputs are used for the transfer \emph{times}, so time and energy are
consistent.

A value of 80~J/MB found in a survey on wireless geophone networks was
discarded: at a typical upload power of 2--4~W it implies about 0.3~Mbit/s,
an IoT modem rather than a 4G smartphone.

\paragraph{Radio session: promotion and tail.} The 4G radio of a phone is
managed by a state machine (RRC, Radio Resource Control) with two main states:
\begin{itemize}
  \item \textbf{IDLE}: the radio is almost off and only wakes up now and then
        to check whether the network is looking for it. Very low power, but
        the phone \emph{cannot transmit}.
  \item \textbf{CONNECTED}: the cell has assigned radio resources to the
        phone, which can send and receive data. Much higher power.
\end{itemize}
\emph{Promotion} is the transition IDLE $\to$ CONNECTED: before sending or
receiving anything, the phone has to negotiate radio resources with the cell.
It costs time and energy even if a single byte follows. \emph{Tail} is what
happens after the last packet: the radio does \emph{not} go back to IDLE at
once, but stays CONNECTED for a fixed inactivity timer, in case more data
arrives (this avoids paying another promotion). During the tail the phone
transmits nothing but still consumes energy; modern modems sleep most of the
time and wake up only periodically (discontinuous reception, DRX).
\[
\text{IDLE}\ \xrightarrow{\ \text{promotion}\ }\ \text{CONNECTED (data)}\
\xrightarrow{\ \text{tail}\ }\ \text{IDLE}
\]

\begin{center}
\begin{tabular}{lcc}
\toprule
& Huang et al.\ 2012 (LTE, 2011 modem) & Narayanan et al.\ 2021 (Verizon 4G) \\
\midrule
promotion & \SI{1.21}{\watt} $\times$ \SI{0.260}{\second} & \SI{0.265}{\second} (power not given) \\
tail      & \SI{1.06}{\watt} $\times$ \SI{11.6}{\second}  & \SI{0.178}{\watt} $\times$ \SI{10.2}{\second} (Tab.~2, Tab.~7) \\
one session & \SI{12.6}{\joule} & \SI{2.14}{\joule} \\
\bottomrule
\end{tabular}
\end{center}
We use Narayanan for the tail and the promotion time, and keep Huang for the
promotion power (about \SI{0.3}{\joule}, negligible). Modern modems spend much
less in the tail than the 2012 one.

\paragraph{Why it matters for us.} Promotion and tail give every
\emph{isolated} transfer a fixed cost, independent of its size (in the
literature this is known as \emph{tail energy}):
\begin{itemize}
  \item \textbf{Control messages.} A SAGE-smart proposal is about 1~KB: as
        payload alone it would cost about \SI{0.002}{\joule}, i.e.\ nothing.
        With the session it costs about \SI{2.1}{\joule}. Without this term,
        querying clients would look free and the negotiation would have no
        cost to balance.
  \item \textbf{Model exchange.} Download and upload are \emph{two} sessions
        per round: training sits in between and lasts longer than the tail,
        so the radio goes back to IDLE before the client has finished.
  \item \textbf{Accepted proposal.} It is immediately followed by the model
        download, so it travels in the same session and pays only its
        payload. This is why in \texttt{fisica2}, where nobody refused, the
        whole negotiation cost \SI{0.001}{\watt\hour} per run.
\end{itemize}

\paragraph{Effect.}
\begin{center}
\begin{tabular}{lrr}
\toprule
& before & now \\
\midrule
model exchange (1.08 MB down + up, two sessions) & \SI{29.2}{\joule} & \SI{8.5}{\joule} \\
\quad of which payload & \SI{4.0}{\joule} & \SI{4.3}{\joule} \\
\quad of which sessions & \SI{25.2}{\joule} & \SI{4.3}{\joule} \\
isolated control message (1 KB) & \SI{12.6}{\joule} & \SI{2.1}{\joule} \\
\bottomrule
\end{tabular}
\end{center}
Communication is now about 3 times cheaper, and the negotiation of SAGE-smart
(control messages) costs less.

%======================================================================
\section{Recharge}
\label{sec:recharge}

\paragraph{Per hour, not per round.} Before, a device started charging with
probability $0.005\,(1-\mathrm{SoC})$ \emph{per round}: an algorithm with
longer rounds saw less charging per hour of real time. Now it is a rate per
hour, applied to the duration $\Delta t$ (hours) of the round:
\[
p = 1 - \exp\!\big(-\lambda\,(1-\mathrm{SoC})\,\Delta t\big).
\]

\paragraph{Source for $\lambda$.} Ferreira et al.\ (Pervasive 2011) followed
the battery of 4035 Android users for four weeks:
\begin{itemize}
  \item charging starts mostly at about 40\,\% (6--8 p.m.) and about 30\,\%
        (1--2 a.m.), plus top-ups at about 80\,\% (8 a.m.);
  \item the daily minimum is 30\,\% on average: users avoid going lower;
  \item the average battery level over the day is 67\,\% (65--74\,\%);
  \item a full charge takes 1--2 hours; sessions are bimodal (many under
        30 minutes, many overnight), 3\,h 54\,min on average.
\end{itemize}
The paper does not give the number of charges per day. From the daily
schedule we infer about \textbf{one main charge per day} per user. With the
average SoC of 0.67 our rate gives one charge per day when
$\lambda\,(1-0.67) = 1/24$, i.e.\ $\lambda = 0.126$/h. This is also the value
that reproduced the old calibration ($0.005$ per round with rounds of about
140~s gives 0.129/h), and it keeps the share of devices charging below the
5\,\% indicated by the supervisor.

\paragraph{Limit: our users are pessimistic.} Real users start charging at
30--40\,\% and almost never go below 30\,\%; ours let the phone reach zero.
Modelling Ferreira's users faithfully would leave almost no depleted device
with any algorithm. Our scenario is therefore a worst case: users that do not
react to a low battery.

%======================================================================
\section{Batch size}
\label{sec:batch}

\paragraph{Why a sweep.} With constant power the physics no longer has an
optimal batch: the energy per sample always decreases with $B$. But a smaller
batch makes more SGD steps per epoch, which can improve accuracy (and, with
non-IID data, increase client drift). The simulator cannot predict which
effect wins, because the training is real. We measured it: FedAvg (linear
battery) and SAGE-smart (datasheet battery) with a \emph{fixed} $B$, three
$\beta$, one seed (the 05/10 report showed that the battery world moves the
results much less than the algorithm).
Accuracy is the mean of the last 20 rounds, SoC the final mean SoC.

\begin{table}[H]
\centering
\caption{Batch sweep, mean over the three $\beta$ (accuracy / final SoC /
depleted devices / Wh spent).}
\label{tab:sweep}
\begin{tabular}{rcc}
\toprule
$B$ & FedAvg & SAGE-smart \\
\midrule
8   & 0.678 / 0.312 / 12.0 / 294 & \textbf{0.702} / 0.500 / 0 / 132 \\
16  & 0.651 / 0.353 / 4.3 / 242  & 0.667 / 0.543 / 0 / 102 \\
32  & 0.631 / 0.376 / 3.3 / 209  & 0.632 / 0.575 / 0 / 87 \\
64  & 0.600 / 0.430 / 0.7 / 191  & 0.608 / 0.574 / 0 / 80 \\
128 & 0.576 / 0.457 / 0.7 / 180  & 0.586 / \textbf{0.578} / 0 / 77 \\
256 & ---                        & 0.531 / 0.576 / 0 / 78 \\
\bottomrule
\end{tabular}
\end{table}

\begin{itemize}
  \item \textbf{A smaller batch always gives more accuracy}, for both
        algorithms and all three $\beta$: FedAvg goes from 0.576 to 0.678,
        SAGE-smart from 0.531 to 0.702.
  \item \textbf{At the same $B$ the accuracy of the two algorithms is about the
        same} (difference 0--2.4 points). The difference is in the battery:
        FedAvg always does 5 epochs and pays the small batches with depleted
        devices (12 at $B=8$, 16 at $\beta=0.1$); SAGE-smart adapts the number
        of epochs and keeps 0 depleted devices, a final SoC of 0.50--0.58 and
        less than half the energy.
  \item On the accuracy--SoC plane, SAGE-smart at $B=8$ dominates
        \emph{every} FedAvg configuration. \emph{SAGE-smart can afford small
        batches because it adapts $E$; FedAvg cannot.}
  \item \textbf{Ablation for the thesis}: part of the accuracy advantage of
        SAGE-smart over FedAvg at $B=128$ comes from the batch, not from the
        negotiation. The FedAvg runs at small $B$ measure how much.
\end{itemize}

\paragraph{The choice.} The baselines stay at $B=128$ (as in the SAGE paper);
only SAGE-smart chooses $B$. For SAGE-smart the SoC is flat from 32 upwards
and drops below 32. $B=16$ is the point closest to the ideal (maximum accuracy
and maximum SoC together): distance 0.050, against 0.078 for $B=8$ and 0.070
for $B=32$.

\paragraph{The rule.} The server proposes $B\in\{16, 32, 64\}$ in the envelope.
The client picks the batch that is \emph{cheapest per SGD step}: it keeps the
smallest $B$ and reduces the epochs first, and moves to the next $B$ only if
not even $E_{lo}=2$ epochs fit in its battery or in the deadline. In practice
a fast or charged client trains at $B=16$, a client with a large partition or
a busy phone goes to 32 or 64. The epoch cap stays at 10, now written as twice
the standard work ($2\times E$ of the baselines).

%======================================================================
\section{A bug found on the way: FedNova with different batches}

SAGE-smart aggregates with FedNova: the update of client $i$ is divided by
its number of steps $\tau_i$ and multiplied by
$\tau_\text{eff}=\sum_i p_i \tau_i$. A client that made few steps is amplified
by $\tau_\text{eff}/\tau_i$. With $B$ ranging from 16 to 256, a client that
missed the deadline fell back to $B=256$ and made about 18 steps against
about 290 of the others (spread 15--16). In the 3-round test the model
collapsed at the first round: accuracy 0.10, loss 9--11.

Two causes and two fixes:
\begin{enumerate}
  \item The amplification was also applied to the \textbf{BatchNorm
        statistics} (running mean and variance) of ResNet-20, which are not
        gradient steps: they were extrapolated beyond the values of any
        client. Now they are averaged as in FedAvg (weights $p_i$), and
        FedNova is applied only to the trainable parameters.
  \item The \textbf{batch range is capped at 64}: the fallback client now
        makes about 70 steps. Over all combinations of tier, partition size,
        battery and workload, the worst-case spread goes from 22.5 to 5.8.
\end{enumerate}
After the fix the same test gives accuracy 0.20 $\to$ 0.29 $\to$ 0.34, mean
$B$ 40 / 32 / 28, spread 4.1--4.4. The quick test now also fails if the
spread exceeds 8. The sweep is not affected: there $B$ was the same for all
clients.

%======================================================================
\section{Results}
\label{sec:results}

\emph{To be added.} Campaign \texttt{rechargeONdeviceON\_ResNet20\_fisica2}:
11 labels (FedAvg, FedProx, SAGE and ESCS-SD in the three battery worlds,
SAGE-smart in the three battery worlds) $\times$ 3 $\beta$, one seed, 250
rounds, 30 clients, 6 per round.

What we expect from the sweep: far fewer depleted devices than in the 05/10
campaign (FedAvg at $B=128$: 0.7 instead of about 15), because training is
2.4--3.4 times cheaper on the low and mid tiers and communication 3 times
cheaper. Whether to restore some pressure on the battery (more rounds, lower
initial SoC) is to be decided after seeing the full campaign.

\subsection{SAGE-smart: how the negotiation went}

Counts summed over the 250 rounds of each run. \emph{Proposals}: envelopes
sent to candidates (all waves). \emph{Refused}: proposals refused or not
answered. \emph{2nd waves}: rounds in which refused candidates had to be
replaced. \emph{Opt-outs}: clients that accepted but withdrew at training time.
\emph{Late}: participations that missed the deadline (in \% of the
participations).

\begin{table}[H]
\centering
\caption{Negotiation of SAGE-smart, campaign \texttt{fisica2}, seed 2026.}
\label{tab:negotiation}
\begin{tabular}{llrrrrr}
\toprule
$\beta$ & battery & proposals & refused & 2nd waves & opt-outs & late \\
\midrule
\multirow{3}{*}{0.1} & linear    & 1490 & 0 & 0 & 0 & 424 (28.5\,\%) \\
                     & Peukert   & 1490 & 0 & 0 & 0 & 423 (28.4\,\%) \\
                     & datasheet & 1490 & 0 & 0 & 0 & 417 (28.0\,\%) \\
\midrule
\multirow{3}{*}{0.5} & linear    & 1490 & 0 & 0 & 0 & 102 (6.8\,\%) \\
                     & Peukert   & 1490 & 0 & 0 & 0 & 102 (6.8\,\%) \\
                     & datasheet & 1490 & 0 & 0 & 0 &  99 (6.6\,\%) \\
\midrule
\multirow{3}{*}{1.0} & linear    & 1490 & 0 & 0 & 0 &  14 (0.9\,\%) \\
                     & Peukert   & 1490 & 0 & 0 & 0 &  12 (0.8\,\%) \\
                     & datasheet & 1490 & 0 & 0 & 0 &  14 (0.9\,\%) \\
\bottomrule
\end{tabular}
\end{table}

\begin{itemize}
  \item \textbf{No client ever refused.} 1490 proposals are exactly the
        selected clients ($2+3+4+5$ in the first rounds, then 6 per round):
        every proposal was accepted at the first wave, and nobody withdrew at
        training time. With the new physics no battery gets close to the
        \SI{20}{\percent} reserve (the most depleted device of SAGE-smart ends
        above 0.34), so the minimum work is always affordable. The
        accept/refuse mechanism is never exercised in this scenario.
  \item \textbf{The negotiation is almost free}: an accepted proposal travels
        in the same radio session as the model download that follows, so it
        costs only its payload. The control energy of a whole run is
        \SI{0.001}{\watt\hour}; the added latency is \SI{0.335}{\second} per
        round, \SI{84}{\second} over the run.
  \item \textbf{Deadline misses depend on $\beta$}: 28\,\% of the
        participations at $\beta=0.1$, 7\,\% at 0.5, 1\,\% at 1.0. At low
        $\beta$ the Dirichlet partitions are much more unequal in size, and
        the clients with the largest partitions cannot fit two epochs in the
        deadline even at $B=64$: they train anyway and declare the delay.
  \item The battery world does not change the negotiation: the three
        batteries give the same counts within a few late participations.
\end{itemize}

%======================================================================
\section{Limits}

\begin{itemize}
  \item \textbf{One seed}, as before.
  \item \textbf{Phones}: one phone per tier, 2020--21, mid and high end only;
        VGG-11 instead of ResNet-20 for the speed; no thermal throttling.
  \item \textbf{4G}: one phone and one operator for the power, Italian speeds
        from 2018.
  \item \textbf{Users}: Ferreira's study is from 2011 (Android 1.6--2.2,
        1500~mAh batteries, users who installed a battery app). Our users let
        the phone reach zero (worst case).
  \item \textbf{Rounds back to back}: as in the SAGE and ESCS papers, rounds
        follow each other with no pause, and the server's information is
        never more than a few minutes old. A scenario with pauses between
        rounds (and really stale information) would make the client-driven
        decision more valuable, but changes the energy balance and the
        information model of all algorithms. Postponed.
  \item \textbf{Differences from the papers to declare}: SAGE without the
        malfunctions (temporary and permanent drop-outs) and without the stop
        at target accuracy 0.4 + 20 rounds; ESCS with a fixed network quality
        per client and no failed transmissions (the paper redraws it every
        round).
\end{itemize}

%======================================================================
\section{Next steps: which part of SAGE-smart brings the advantage?}
\label{sec:next}

\emph{Reply to Alessio's e-mail}: below is a proposal on how to isolate the
contributions of SAGE-smart, to discuss in the next call.

\paragraph{What has to be explained.} Against FedAvg ($B=128$, $E=5$),
averaged over the three $\beta$, SAGE-smart gains about +5 accuracy points,
uses about half the energy, ends with a mean SoC about 0.1 higher and keeps
its most depleted device above 0.34. SAGE-smart is the SAGE ablation
($a\cdot\mathrm{SoC} + b\cdot D$) plus nine additions. Each can be switched off
on its own:

\begin{table}[H]
\centering
\small
\caption{Components of SAGE-smart that can be isolated. ``Config'' = one
value in \texttt{experiment.toml}; ``code'' = a small switch to add.}
\label{tab:ablations}
\begin{tabular}{p{0.3cm}p{3.1cm}p{4.2cm}p{3.0cm}p{4.0cm}}
\toprule
& component & what it does & how to switch it off & evidence so far \\
\midrule
1 & staleness term $c$ in the score & gives priority to clients that have not trained for a while: spreads participation & config: \texttt{smart\_stale\_weight = 0} & earlier measurement: without it 4/30 clients were never selected and one 228 times out of 250 (Gini 0.62 vs 0.19 of FedAvg); SAGE, which has no such term, drops to 0.30 at $\beta=0.1$ (datasheet) \\
2 & server-side SoC estimate (option 2) & ranks with the SoC extrapolated from the last report & code: last report as is, or the true SoC (upper bound) & rounds back to back: the estimate is almost exact, small effect expected \\
3 & negotiation (option 3) & proposal, accept / refuse, replacement of refusals & config: \texttt{smart\_max\_waves = 1}; code: no proposal at all & 0 refusals in \texttt{fisica2}: no effect expected in this scenario \\
4 & adaptive local epochs & each client chooses $E\in[2,10]$ from the loss of the global model on its data ($u$) and from its battery ($h$) & code: $E=5$ fixed as the baselines; variants with only $u$ or only $h$ & at $\beta=0.1$ SAGE-smart trains 3.4 epochs on average instead of 5: likely the main source of the energy saving \\
5 & late binding (the client decides) & $(E,B)$ chosen by the client with its \emph{current} state instead of by the server with its estimate & code: server-driven variant, the client only executes & core of the thesis; small difference expected while the server's information is fresh \\
6 & smaller batch & $B\in\{16,32,64\}$ instead of 128 & config: \texttt{smart\_batch\_min = smart\_batch\_max = 128} & sweep: FedAvg at $B=16$ gains 7.5 points over $B=128$: part of the accuracy advantage comes from $B$ \\
7 & deadline & caps the round time and the straggler & config: \texttt{smart\_deadline\_mult} very large; variant: deadline at $B=16$ & at $\beta=0.1$ the 38~s deadline pushes clients from $B=16$ to $B\approx 31$: 0.48 instead of 0.58 \\
8 & reserve and energy guard & the client never goes below 20\,\%, stops early if needed & config: \texttt{smart\_soc\_min = 0} & never binding now: no device gets close to 20\,\% \\
9 & FedNova aggregation & normalises each update by its number of steps (BatchNorm averaged) & code: plain data-weighted average & needed with different $(E,B)$: the model collapsed with a step spread of 15--16 \\
\bottomrule
\end{tabular}
\end{table}

\paragraph{Proposed plan.}
\begin{enumerate}
  \item \textbf{Leave-one-out}: full SAGE-smart minus one component at a time,
        datasheet battery only (the battery world changes the results by about
        1.5 points), three $\beta$. Nine variants $\times$ 3 $\beta$ = 27 runs,
        about 18 hours per seed on our single GPU. Components 1, 3, 6, 7, 8 need
        only a configuration value; 2, 4, 5, 9 need a small switch in the code.
  \item \textbf{Cumulative ladder} for the thesis narrative: SAGE ablation
        $\to$ + staleness $\to$ + adaptive $E$ decided by the server $\to$ +
        late binding $\to$ + negotiation $\to$ + small $B$. The two ends already
        exist (SAGE ablation and full SAGE-smart in \texttt{fisica2}); the
        batch step is partly covered by the sweep.
  \item \textbf{Scenario under pressure.} Components 3 (negotiation) and 8
        (reserve) can show their value only when batteries are low; in the
        current scenario they are never used. They should be measured also in
        a ``stress'' scenario: lower initial SoC, or a pause between rounds.
        A first estimate with the physics alone: with 2--5 minutes between
        rounds FedAvg causes about 6--7 more depleted devices than a run with
        no training at all, so a no-training reference run is needed to
        separate the effect of FL from the phone's own use.
\end{enumerate}

\paragraph{Other decisions for the call.}
\begin{itemize}
  \item \textbf{Seeds}: run 2026, 7 and 13, so that differences of a few
        points can be tested.
  \item \textbf{Battery metric}: with the new physics almost no device is
        depleted. We propose to count the devices that end \emph{below
        \SI{20}{\percent}} (battery-saver level). A first estimate gives 7 for
        FedAvg against 0.3 with no training. The new figure (fig10) is ready,
        but the column is not in the \texttt{fisica2} runs: it needs new runs.
  \item \textbf{Deadline of SAGE-smart}: computed on the baselines' batch
        (128, now) or on SAGE-smart's preferred batch (16, longer deadline,
        about +0.1 accuracy at $\beta=0.1$ in the sweep)?
\end{itemize}

%======================================================================
\section*{References}

\begin{itemize}
  \item D.~Xu et al., ``Mandheling: Mixed-Precision On-Device DNN Training
        with DSP Offloading'', ACM MobiCom 2022 (arXiv 2206.07509).
  \item A.~Narayanan et al., ``A Variegated Look at 5G in the Wild:
        Performance, Power, and QoE Implications'', ACM SIGCOMM 2021; dataset
        \url{https://github.com/SIGCOMM21-5G/artifact}.
  \item Opensignal, ``State of Mobile Networks: Italy'', May 2018.
  \item J.~Huang et al., ``A Close Examination of Performance and Power
        Characteristics of 4G LTE Networks'', ACM MobiSys 2012.
  \item D.~Ferreira, A.~K.~Dey, V.~Kostakos, ``Understanding Human-Smartphone
        Concerns: A Study of Battery Life'', Pervasive 2011, LNCS 6696,
        pp.~19--33.
  \item F.~Maciel et al., ``Federated learning energy saving through client
        selection'' (ESCS), Pervasive and Mobile Computing 103 (2024) 101948.
  \item M.~Savoia et al., ``SAGE: Sustainable and green energy-aware client
        selection for federated learning'', Applied Energy 410 (2026) 127568.
  \item J.~Wang et al., ``Tackling the Objective Inconsistency Problem in
        Heterogeneous Federated Optimization'' (FedNova), NeurIPS 2020.
  \item TrainDeeploy (on-device training on GAP9) --- full reference to be
        added.
  \item ``Accurate Battery Model'', datasheet document provided by the
        supervisor.
\end{itemize}

\end{document}
